% !TEX root = CV_Tramarin.tex
% !TEX TS-program = XeLaTeX

\subsectioniten{Dottorato di Ricerca}{PhD}

\cventryit{2009 -- 2011}{Dottorato di Ricerca in Scienza e Tecnologie dell’Informazione, 
	XXIV ciclo}{\newline Scuola di Dottorato in Ingegneria dell’Informazione, Dip. Ingegneria dell’Informazione}{}{}{Università degli Studi di Padova, Italia}
	
    \cvdoubleitemit{Titolo della tesi}{\emph{Industrial Wireless Sensor Networks. Simulation and measurements in an interfering environment}}{Giudizio finale}{``Ottimo''}
    \cvdoubleitemit{Data di difesa}{19 aprile 2012}{Advisor}{Prof. M. Bertocco}
	%\cvitem{SC/SSD}{09/E4 --- ING/INF-07}
    
	
%	\cvitem{}{\textbf{\itshape Scuole di Dottorato frequentate}}
%	\cvlistitem{GMEE Seminario di Eccellenza, Italo Gorini\newline
%		      5 – 9 settembre, 2011, Siena}
%	\cvlistitem{GMEE Seminario di Eccellenza, Italo Gorini\newline
%		      30 agosto – 3 settembre, 2010, Pistoia}
%	\cvlistitem{GMEE Seminario di Eccellenza, Italo Gorini\newline
%		      31 agosto – 4 settembre, 2009, Perugia}

\cventryen{2009 -- 2011}{PhD in Information Science and Technologies, 
    cycle XXIV}{}
    {}{Department of Information Engineering}{University of Padova, Padova, Italy}
    
    \cvdoubleitemen{Thesis title}{\emph{Industrial Wireless Sensor Networks. Simulation and measurements in an interfering environment}}{Final grade}{Excellent (``Ottimo'')}
    \cvdoubleitemen{Defense date}{April 19th, 2012}{Advisor}{Prof. M. Bertocco}
    
%	\begin{minipage}{\textwidth}
%	\cvitem{}{\textbf{\itshape Attended PhD courses (with final grade)}}
%	\cvlistitem{Stochastic Modeling of Computer and  Communication Systems (Prof. A. Willig)}
%	\cvlistitem{Wireless Sensor Networks (Prof. A. Willig)}
%	\cvlistitem{Applied Functional Analysis (Prof. P. Ciatti,  Prof. G. Pillonetto)}
%	\cvlistitem{Statistical Methods (Prof. L. Finesso)}
%	\end{minipage}
%    \fi

\subsectioniten{Laurea}{Master's Degree}

\cventryit{2006 -- 2008}{Laurea Magistrale in Ingegneria Elettronica}
{Dip. Ingegneria dell’Informazione, Università degli Studi di Padova}
{Padova}{}{}

\cvdoubleitemit{Titolo della tesi}{\emph{Cross-Layer Analysis of Interfering Wireless Sensors Networks}}
{Classe}{32/S - Classe delle lauree specialistiche in ingegneria elettronica}

\cvdoubleitemit{Voto di laurea}{110/110 e lode, media dei voti degli esami 29.77/30 su 13 esami}{Relatore}{Prof. M. Bertocco}

%\cvitem{Abilitazione}{2009, Abilitazione Nazionale alla Professione di Ingegnere Elettronico}

% \vspace{5mm}
% \cventry{2003--2006}{Laurea Triennale in Ingegneria dell'Informazione}
% {Dip. Ingegneria dell’Informazione, Università degli Studi di Padova}
% {Padova}{}{}
% \cvitem{Classe}{9 - Classe delle lauree in ingegneria dell'informazione}
% \cvitem{Titolo della tesi}{\emph{Studio e realizzazione del decodificatore di 
% Viterbi per il sistema UWB di tipo MB-OFDM nella modalità trasmissiva DCM}}
% \cvitem{Relatore}{Prof. N. Laurenti}
% \cvitem{Voto di laurea}{110/110 e lode, media dei voti degli esami 29.62/30 su 23 esami}


\cventryen{2006--2008}{Master's Degree in Electronic Engineering}
{Department of Information Engineering, University of Padova}
{Padova}{Italy}{}
%\cvdoubleitem{Thesis}{\emph{Cross-Layer Analysis of Interfering Wireless Sensors Networks}}{Advisor}{Prof. M. Bertocco}

\cvdoubleitemen{Thesis title}{\emph{Cross-Layer Analysis of Interfering Wireless Sensors Networks}}
{Supervisor}{Prof. M. Bertocco}
\cvitemen{Grade}{110/110 \textit{(cum laude)}, GPA 29.77/30 over 13 exams}

%\cvdoubleitemit{Voto di laurea}{110/110 e lode, media dei voti degli esami 29.77/30 su 13 esami}